% Phase 3 CP3b. Modeling line: does the learned operator predict? Table 1 = tables/tab_prediction.tex.
\subsection{Predictive Accuracy of the Learned Operator}
\label{sec:exp_prediction}

\textbf{Question.} RQ1a and RQ1b: whether multi-turn history improves prediction over the current reading alone, and whether a learned nonlinear lifting or a recurrent sequence model beats the linear special case reported as ours.
\textbf{Setup.} Setup, tasks and baselines follow \S\ref{sec:exp_setup}; Table~\ref{tab:prediction} reports the skill score defined there.

% GENERATED by tools/build_tables.py. Do not edit by hand; edit the evidence ledger instead.
% MAIN-TEXT TABLE 1 of 2 (modeling line). D18: the main text carries exactly two tables.
\begin{table}[t]
\centering
\caption{Free-running prediction skill on seven multi-turn tasks. Higher is better. Best per column in bold, second underlined.}
\label{tab:prediction}
\footnotesize
\resizebox{\textwidth}{!}{
\begin{tabular}{lccccccc}
\toprule
Predictor & Sentence len. & Items (s2) & Items (s1) & Char. len. & Formality & Constraint & CEFR \\
\midrule
Stateless & 0.000 & 0.000 & 0.000 & 0.000 & 0.000 & 0.000 & 0.000 \\  % n_main_002 n_main_034 n_main_018 n_main_066 n_main_082 n_main_121 n_main_169
Last-turn & 0.166 & 0.514 & 0.395 & 0.034 & \textbf{0.414} & 0.285 & 0.518 \\  % n_main_006 n_main_038 n_main_022 n_main_070 n_main_086 n_main_123 n_main_171
Memory w/o action & \underline{0.764} & \underline{0.824} & 0.615 & \underline{0.291} & 0.357 & 0.517 & \underline{0.561} \\  % n_main_008 n_main_040 n_main_024 n_main_072 n_main_088 n_main_125 n_main_173
LSTM & 0.750 & 0.808 & \underline{0.634} & \textbf{0.301} & 0.373 & \textbf{0.546} & -0.101 \\  % n_main_012 n_main_044 n_main_028 n_main_076 n_main_092 n_main_127 n_main_175
KOMPAS (learned lifting) & 0.752 & 0.821 & 0.627 & 0.239 & \underline{0.385} & 0.471 & -0.032 \\  % n_main_010 n_main_042 n_main_026 n_main_074 n_main_090 n_main_129 n_main_177
KOMPAS (linear) & \textbf{0.769} & \textbf{0.826} & \textbf{0.640} & 0.269 & 0.379 & \underline{0.519} & \textbf{0.567} \\  % n_main_004 n_main_036 n_main_020 n_main_068 n_main_084 n_main_131 n_main_179
\bottomrule
\end{tabular}
}
\end{table}


\textbf{Findings, ranking and history.} The linear special case, KOMPAS (linear), ranks first on four of seven columns: sentence length, Joint-3, Joint-2 and CEFR.  % c4; ours_best_of_7=4 (n_main_004 n_main_036 n_main_020 n_main_179)
It ranks second on the remaining three: sentiment, defense and constraint.
The driver is that all three second places are close: the linear case trails the top method by no more than 0.034 in point-estimate skill score.  % c4 driver_keyword  % ours_second_of_7=3 (n_main_100 sentiment; n_main_147 defense; n_main_131 constraint)
History beyond the current turn beats the current reading on four of seven columns, with paired bootstrap intervals excluding zero.  % c1; contrasts.ours_minus_last_turn (n_main_013 n_main_045 n_main_029 n_main_109 n_main_148 n_main_132 n_main_180)
The three exceptions are Joint-2 and sentiment, whose intervals contain zero, and defense, whose point estimate is negative with an interval containing zero.  % n_main_029 Joint-2 (vector_count_stage1); n_main_109 sentiment; n_main_148 defense
The driver is that the window exposes both level and trend, which a single reading cannot carry.  % m1

\textbf{Findings, lifting and LSTM.} A learned nonlinear lifting never beats the identity lifting with an interval excluding zero, on any of the seven columns.  % c2; contrasts.ours_minus_nonlinear
It loses on three of them, with the largest gap on CEFR at +0.599 in favor of the linear case.  % contrasts.ours_minus_nonlinear:tsar_cefr=0.599 (n_main_183)
The driver is that the attribute readouts are low-dimensional, one scalar per attribute and at most three per task, leaving a learned lifting little structure to expose.  % m4; Joint-2/Joint-3 are 2- and 3-dimensional, the rest scalar
Against an LSTM sequence baseline, comparisons split three ways with intervals excluding zero: two favor the linear case, one favors the LSTM.  % c3; contrasts.ours_minus_lstm
Joint-3 and CEFR favor the linear case, while constraint favors the LSTM.  % n_main_048 n_main_182 favor linear; n_main_134 favor LSTM
No mechanism is established for the LSTM win, and none is proposed here.  % c3 driver_keyword: no mechanism established

Together, \uline{history beyond the current turn beats the current reading on four of seven tracked attributes, which answers RQ1a yes on those four. Neither a learned lifting nor an LSTM sequence model beats the reported linear special case on balance, which answers RQ1b no.}
